\documentclass[11pt]{article}
\usepackage[T1]{fontenc}
\usepackage{amsmath,amssymb,amsthm}
\usepackage[margin=1in]{geometry}
\usepackage{hyperref}
\newtheorem{theorem}{Theorem}
\newtheorem{lemma}[theorem]{Lemma}
\newtheorem{corollary}[theorem]{Corollary}
\newcommand{\ddc}{\frac{i}{2\pi}\partial\bar\partial}
\newcommand{\tr}{\operatorname{tr}}
\title{A spectral gap forces collapse of normalized Gauduchon tests}
\author{}
\date{11 October 2026}
\begin{document}
\maketitle
\noindent\textbf{Research note.}
The restricted results passed independent review and statement-only
reconstruction. The original P09 problem is not settled by this note.

\begin{abstract}
We bound the determinant denominator of every normalized Hermitian test
by the pairing with a semipositive form having a fixed spectral gap.
If that form has rank \(q<n\), a pairing of size \(e\) forces the denominator
to be \(O(e^{q/n})\). On a product of elliptic curves the estimate is sharp,
with an explicit maximizing metric at every prescribed pairing.
A closed strictly positive background on this product has positive root
volume and positive strict-positivity threshold, yet the determinant
denominator tends to zero along every sequence minimizing that threshold.
This rules out a universal lower bound for that particular normalization
of threshold-minimizing tests. The general positive-intersection
root-volume problem is not settled here.
\end{abstract}

\section{The two quotients}

Let \(X\) be a connected compact complex manifold of dimension \(n\geq2\),
\(h\) a smooth positive Hermitian form, and \(\mu\) a smooth positive
probability volume form. A smooth positive Hermitian form \(\gamma\) is
Gauduchon when \(\ddc(\gamma^{n-1})=0\). Put
\[
 A_\theta(\gamma)=\int_X\theta\wedge\gamma^{n-1},
 \qquad
 D_\mu(\gamma)=\int_X
       \left(\frac{\mu}{\gamma^n}\right)^{1/n}\gamma^n .
\]
For a smooth positive background \(\alpha\), the two quantities considered
here are
\[
 \kappa_h(\alpha)=\inf_{\gamma\ {\rm Gauduchon}}
       \frac{A_\alpha(\gamma)}{A_h(\gamma)},
 \qquad
 \lambda_\mu(\alpha)=\inf_{\gamma\ {\rm Gauduchon}}
       \frac{A_\alpha(\gamma)}{D_\mu(\gamma)}.
\]
The factor \(n\) in the usual trace identity is included in these
definitions. Multiplication of \(\gamma\) by \(t>0\) multiplies
all three pairings \(A_\alpha,A_h,D_\mu\) by \(t^{n-1}\).
Thus each ray can be normalized by \(A_h(\gamma)=1\).

The source definitions and the proposed threshold-sequence argument
are recorded in \cite{volume,collapse,bridge}. They are used as definitions
and as the formulation to be tested. All estimates below are proved
directly, for arbitrary smooth positive Hermitian tests; their validity
does not depend on a solution theorem or on a restriction to invariant
Gauduchon tests.

\section{A spectral estimate for every test}

Write
\[
 H=\int_Xh^n,\qquad \nu=\frac{h^n}{H},\qquad
 r=\frac{\mu}{\nu},\qquad R=\max_X r.
\]
The maximum \(R\) is finite and positive. At a point, an \(h\)-unitary
frame identifies \(\gamma\) with a positive Hermitian matrix \(G\).
Set \(C=(\det G)G^{-1}\). This is the adjugate endomorphism; the convention
is independent of the unitary frame. It satisfies
\[
 \det C=(\det G)^{n-1},\qquad
 A_h(\gamma)=\frac Hn\int_X\tr C\,d\nu,\qquad
 D_\mu(\gamma)=H^{(n-1)/n}
       \int_X r^{1/n}(\det C)^{1/n}\,d\nu .
\]
These identities can be checked in a unitary frame diagonalizing \(G\):
\(\gamma^n=(\det G)h^n\), while
\(h\wedge\gamma^{n-1}/h^n=\tr C/n\).
The same expansion gives
\(B\wedge\gamma^{n-1}/h^n=\tr(BC)/n\), including an arbitrary
Hermitian \(B\). Substituting \(\mu=rh^n/H\) gives the displayed
denominator formula. Unitary invariance then gives the identities globally.

\begin{theorem}\label{thm:gap}
Let \(B\) be a smooth semipositive real \((1,1)\)-form of constant
complex rank \(q\), where \(1\leq q<n\), and put \(m=n-q\).
Suppose its positive eigenvalues relative to \(h\) are at least
\(\delta>0\). For every smooth positive Hermitian \(\gamma\) with
\(A_h(\gamma)=1\), put \(e=A_B(\gamma)\).
If \(0<e\leq\delta q/n\), then
\begin{equation}\label{eq:gap}
 D_\mu(\gamma)\leq
 n(R/H)^{1/n}
 \left(\frac{1-e/\delta}{m}\right)^{m/n}
 \left(\frac{e/\delta}{q}\right)^{q/n}.
\end{equation}
In particular \(D_\mu(\gamma)\leq
n(R/H)^{1/n}m^{-m/n}q^{-q/n}\delta^{-q/n}e^{q/n}\).
The estimate includes variable and nondiagonal tests.
\end{theorem}

\begin{proof}
Let \(P\) be the \(h\)-orthogonal projection onto the positive spectral
subbundle of \(B\), and \(Q=I-P\). Constant rank and the positive gap make
these smooth subbundles. In matrix notation \(B\geq\delta P\).
Put \(p=\tr(PC)>0\) and \(s=\tr(QC)>0\).
In the orthogonal splitting of dimensions \(m,q\), write
\[
 C=\begin{pmatrix}U&V\\V^*&W\end{pmatrix}.
\]
The blocks \(U,W\) are positive. The Schur complement
\(W-V^*U^{-1}V\) is positive and at most \(W\), so
\[
 \det C
 =\det U\,\det(W-V^*U^{-1}V)
 \leq\det U\,\det W
 \leq (s/m)^m(p/q)^q.
\]
For completeness, the determinant inequality follows by conjugating
the Schur complement by \(W^{-1/2}\): its eigenvalues belong to \((0,1]\).
The last step is the arithmetic--geometric mean inequality for the
positive eigenvalues of each diagonal block.

Let \(z=\int_Xp\,d\nu\) and \(y=Hz/n\). Normalization gives
\[
 \int_X(s+p)\,d\nu=n/H,\qquad
 \int_Xs\,d\nu=n/H-z.
\]
Also
\[
 e=\frac Hn\int_X\tr(BC)\,d\nu\geq\frac{H\delta}{n}z.
\]
Indeed \(\tr((B-\delta P)C)\geq0\), as seen after conjugation by
\(C^{1/2}\). Hence \(0<y<1\) and \(y\leq e/\delta\).

The determinant estimate and weighted arithmetic--geometric mean give
\begin{align*}
 D_\mu(\gamma)
 &\leq H^{(n-1)/n}R^{1/n}
       \left(\frac{\int_Xs\,d\nu}{m}\right)^{m/n}
       \left(\frac{\int_Xp\,d\nu}{q}\right)^{q/n}\\
 &=n(R/H)^{1/n}
       \left(\frac{1-y}{m}\right)^{m/n}
       \left(\frac yq\right)^{q/n}.
\end{align*}
This integral inequality follows by dividing \(s,p\) by their positive
integrals and applying weighted arithmetic--geometric mean pointwise,
then integrating. It therefore imposes no commutativity or constancy
hypothesis on \(C\).

The logarithmic derivative of
\((1-y)^{m/n}y^{q/n}\) is
\((q-ny)/(ny(1-y))\). It is nonnegative for \(0<y\leq q/n\).
Our assumptions put both \(y\) and \(e/\delta\) in this interval.
Replacing \(y\) by \(e/\delta\) proves \eqref{eq:gap}.
Discarding the factor \((1-e/\delta)^{m/n}\leq1\) proves the last bound.
\end{proof}

\begin{corollary}\label{cor:sequence}
Suppose \(a>0\), \(\alpha-a h=B\) satisfies the hypotheses of
Theorem~\ref{thm:gap}, and \(\kappa_h(\alpha)=a\).
Then every normalized sequence of Gauduchon metrics with
\(A_\alpha(\gamma_j)\to a\) satisfies
\[
 D_\mu(\gamma_j)=
 O\bigl((A_\alpha(\gamma_j)-a)^{q/n}\bigr)\longrightarrow0.
\]
\end{corollary}
\begin{proof}
Here \(e_j=A_B(\gamma_j)=A_\alpha(\gamma_j)-a\to0\).
Every \(e_j\) is strictly positive because \(B\) has positive rank and
\(\gamma_j\) is positive. Apply Theorem~\ref{thm:gap} for all sufficiently
large \(j\). The equality \(\kappa_h(\alpha)=a\) is a hypothesis of this
corollary, not an inference from the pointwise lower bound alone.
\end{proof}

\section{A closed positive example and sharp constants}

Let \(E=\mathbb C/(\mathbb Z+i\mathbb Z)\), let \(X=E^n\), and let
\(\omega_i\) be the pullback of the translation-invariant positive
area form on the \(i\)-th factor with integral one. Set
\[
 \mu=\omega_1\wedge\cdots\wedge\omega_n,\qquad
 h=\sum_{i=1}^n\omega_i,\qquad
 B=\delta\sum_{i=m+1}^n\omega_i,\qquad
 \alpha=a h+B,
\]
where \(a,\delta>0\), \(1\leq m<n\), and \(q=n-m\).
The forms \(h,\alpha\) are closed and strictly positive. In particular
\(\alpha\) is nef, is itself a strict positive representative, and
\[
 \int_X\alpha^n=n!a^m(a+\delta)^q>0.
\]

\begin{theorem}\label{thm:flat}
For this fixed geometry,
\[
 \kappa_h(\alpha)=a,\qquad
 \lambda_\mu(\alpha)=\bigl(n!a^m(a+\delta)^q\bigr)^{1/n}>0.
\]
For every \(e\in(0,\delta)\), among all smooth normalized Gauduchon
metrics with \(A_\alpha(\gamma)=a+e\), the maximum denominator is
\begin{equation}\label{eq:flat}
 \frac{n}{(n!)^{1/n}}
 \left(\frac{1-e/\delta}{m}\right)^{m/n}
 \left(\frac{e/\delta}{q}\right)^{q/n}.
\end{equation}
A translation-invariant closed metric attains this value.
Nevertheless every normalized threshold-minimizing sequence has
\(D_\mu(\gamma_j)\to0\).
\end{theorem}

\begin{proof}
In this normalization \(H=n!\), \(\nu=\mu\), and \(R=1\).
Moreover \(B=\delta P\), so the preceding proof gives
\(y=e/\delta\) exactly. Its determinant and integral estimates give
\eqref{eq:flat} for every \(e\in(0,\delta)\); the monotonicity step of
Theorem~\ref{thm:gap} is unnecessary.

We verify attainment within the full test class. Put \(x=e/\delta\), and
choose weights
\[
 w_i=\frac{1-x}{m}\quad(1\leq i\leq m),\qquad
 w_i=\frac{x}{q}\quad(m<i\leq n).
\]
Let \(C\) be the constant diagonal matrix with entries
\(c_i=w_i/(n-1)!\), and set
\[
 G=(\det C)^{1/(n-1)}C^{-1}.
\]
It is positive, and direct determinants give
\(\det G=(\det C)^{1/(n-1)}\); hence \((\det G)G^{-1}=C\).
The associated translation-invariant metric is closed and hence
Gauduchon. Directly,
\[
 A_h=(n-1)!\sum_i c_i=1,\qquad
 A_B=\delta(n-1)!\sum_{i>m}c_i=\delta x=e.
\]
Its denominator is
\[
 (n!)^{(n-1)/n}(\det C)^{1/n}
 =\frac{n}{(n!)^{1/n}}\left(\prod_i w_i\right)^{1/n},
\]
which equals \eqref{eq:flat}.

Since \(\alpha\geq a h\), every test satisfies \(A_\alpha/A_h\geq a\).
The constructed metrics have \(A_\alpha=a+\delta x\), and \(x\downarrow0\)
therefore proves \(\kappa_h(\alpha)=a\).
For an arbitrary normalized threshold-minimizing sequence,
\(e_j=A_\alpha(\gamma_j)-a\to0\).
Either \eqref{eq:flat} or Corollary~\ref{cor:sequence} then proves
denominator collapse, including for variable nondiagonal metrics.

It remains to compute the root volume without a solution theorem.
The pointwise arithmetic--geometric mean inequality for the eigenvalues
of \(\alpha\) relative to \(\gamma\) yields
\[
 \alpha\wedge\gamma^{n-1}
 \geq
 \left(\frac{\alpha^n}{\gamma^n}\right)^{1/n}\gamma^n.
\]
Here \(\alpha^n=M\mu\), where \(M=n!a^m(a+\delta)^q\).
Thus \(A_\alpha(\gamma)\geq M^{1/n}D_\mu(\gamma)\) for every smooth
positive metric, in particular every Gauduchon metric.
The closed metric \(\gamma=\alpha\) gives equality. Taking the infimum
proves the asserted formula for \(\lambda_\mu\).
\end{proof}

\section{What the example decides}

The example has a closed nef background with positive top intersection,
positive root volume, and an explicit strict positive representative.
It has no normalized threshold-minimizing sequence whose determinant
denominators have positive lower limit. Thus the existence of such
a sequence cannot be a necessary consequence of closed nef positive
intersection, even on a compact K\"ahler manifold.

There is no contradiction with the conditional collapse proposition:
its nonbigness hypothesis is absent here, and the hypothesis of a positive
denominator lower limit in its second assertion simply fails.
Nor is there a counterexample to positive-intersection bigness.
The general first bridge on arbitrary compact complex manifolds,
and the classical closed nef bigness conjecture, are not decided.
Any continuation using normalized threshold sequences must distinguish
a useful sufficient condition from a universally available condition.

The two quotient problems are genuinely different in this example.
The metric \(\alpha\) attains the root-volume quotient. After normalization,
its excess over \(a\) is
\[
 e_\lambda=\frac{\delta q a}{na+m\delta}>0,
\]
because its normalized cofactor weights are
\(w_i=a_i^{-1}/\sum_\ell a_\ell^{-1}\), with \(a_i=a\) on the first
block and \(a_i=a+\delta\) on the second. Threshold-minimizing
sequences instead have excess tending to zero and cannot keep a positive
determinant denominator.

This note does not replace the frozen P09 manuscripts or change any
previous native verdict, acceptance record, publication count, or
original-problem completion flag. Literature priority is not asserted.

\begin{thebibliography}{3}
\bibitem{volume}
M. Xia, \emph{Singularities in Hermitian pluripotential theory},
working manuscript: \emph{A Gauduchon variational volume},
source snapshot \texttt{0646d2cf}, definition and dual-formula context.
\href{https://github.com/MingchenXia/non-kahler-pluripotential-theory/blob/0646d2cf9d9bb10e15bc357305cce3c58ecb4741/book/author/chapter_gauduchon_variational_volume.tex}{Versioned source}.

\bibitem{collapse}
Same manuscript, \emph{Collapse of normalized Gauduchon tests},
entire subsection, source snapshot \texttt{0646d2cf}.
\href{https://github.com/MingchenXia/non-kahler-pluripotential-theory/blob/0646d2cf9d9bb10e15bc357305cce3c58ecb4741/book/author/section_gauduchon_collapse.tex}{Versioned source}.

\bibitem{bridge}
Same manuscript, \emph{Closed positivity and geometric equations},
closed-nef bridge problem and its surrounding explanation,
source snapshot \texttt{0646d2cf}.
\href{https://github.com/MingchenXia/non-kahler-pluripotential-theory/blob/0646d2cf9d9bb10e15bc357305cce3c58ecb4741/book/author/section_non_kahler_geometry.tex}{Versioned source}.
\end{thebibliography}
\end{document}
